\documentclass[11pt]{article}
\usepackage[T1]{fontenc}
\usepackage{amsmath,amssymb,amsthm}
\usepackage[margin=1in]{geometry}
\usepackage{graphicx,microtype,xcolor}
\usepackage[hidelinks]{hyperref}
\newtheorem{theorem}{Theorem}
\newcommand{\E}{\mathbb E}
\newcommand{\Pp}{\mathbb P}
\title{A counterexample to the all-vector law of the iterated logarithm for random weighted shifts}
\author{Automated Math Discovery Run \texttt{fa\_banach\_001}}
\date{21 August 2026}
\begin{document}
\maketitle
\begin{center}
\fcolorbox{black}{red!5}{\parbox{0.88\linewidth}{
\textbf{Result.} Conjecture~7.17 of Cheng--Fang--Zhu is false.  There are
bounded positive i.i.d. weights satisfying all of the paper's hypotheses
and a deterministic nonzero vector \(x\in\ell^2\) for which
\[
 \frac{\log\|T^n x\|^2}{\sqrt{2\sigma^2n\log\log n}}\longrightarrow+\infty
 \quad\text{almost surely}.
\]
}}
\end{center}

\section{The source conjecture}
For a unilateral random weighted shift \(Te_k=X_ke_{k+1}\), Lemma~7.15 of
\cite{CFZ} proves, for each basis vector \(e_k\), the usual LIL with
normalization \(\sqrt{2\sigma^2n\log\log n}\), where
\(\E\log X_1=0\) and \(\sigma^2=\E(\log X_1^2)^2<\infty\).
Conjecture~7.17 asks whether the same assertion holds for every
\(x\in\ell^2\).

\begin{figure}[ht]
\centering
\includegraphics[width=0.96\linewidth,height=0.58\textheight,keepaspectratio]{figures/open_problem_crop.png}
\caption{Lemma~7.15, formula (49), and Conjecture~7.17 in arXiv:1811.05761,
PDF page 48.}
\end{figure}

\section{Proof intuition}
For a fixed basis vector, \(\|T^n e_k\|\) sees only one length-\(n\)
weight block, so the ordinary random-walk LIL applies.  An infinite-support
vector is different: its norm is a sum over all starting locations and can
select a rare exceptionally favorable block far out in the tail.

With symmetric weights \(e^{\pm1}\), an all-\(e\) block has probability
\(2^{-n}\).  Among \(\lceil e^n\rceil\) disjoint blocks, the probability of
seeing none is double-exponentially small and hence summable.  The
coefficient \(k^{-3/5}\) decays slowly enough that an all-\(e\) block starting
before \(n e^n\) still contributes \(\exp((4/5)n-O(\log n))\) to the squared
norm.  Borel--Cantelli therefore upgrades this rare-block mechanism to an
almost-sure lower bound for every sufficiently large \(n\), far above the
\(\sqrt{n\log\log n}\) LIL scale.

\section{Counterexample}
\begin{theorem}
Let \(X_j=e\) and \(X_j=e^{-1}\), each with probability \(1/2\), independently,
and set
\[
 x=\sum_{k=1}^{\infty}k^{-3/5}e_k.
\]
Then \(x\in\ell^2\), the weights satisfy the hypotheses of Lemma~7.15 with
\(\sigma^2=4\), and almost surely
\[
 \lim_{n\to\infty}\frac{\log\|T^n x\|^2}{\sqrt{8n\log\log n}}=+\infty.
\]
In particular Conjecture~7.17 is false.
\end{theorem}

\begin{proof}
Plainly \(\sum k^{-6/5}<\infty\), while
\(\E\log X_1=0\) and
\(\E(\log X_1^2)^2=\E(\pm2)^2=4\).

Fix \(n\), put \(M_n=\lceil e^n\rceil\), and consider the \(M_n\) disjoint
blocks of indices
\[
 B_{n,r}=\{1+rn,\ldots,(r+1)n\},\qquad 0\le r<M_n.
\]
Let \(F_n\) be the event that none of these blocks consists entirely of
weights equal to \(e\).  Independence within this fixed collection gives
\[
 \Pp(F_n)=(1-2^{-n})^{M_n}
 \le \exp(-M_n2^{-n})
 \le \exp\bigl(-(e/2)^n\bigr).
\]
The last sequence is summable, so the first Borel--Cantelli lemma implies
that, almost surely, \(F_n\) occurs for only finitely many \(n\).  Notice that
no independence between different values of \(n\) is needed.

Hence, almost surely for every sufficiently large \(n\), there is an
\(r<M_n\) such that all weights in \(B_{n,r}\) equal \(e\).  With
\(k=1+rn\le nM_n\), formula (49) of the source then yields the one-term
lower bound
\[
 \|T^n x\|^2
 =\sum_{j\ge1}j^{-6/5}X_j^2\cdots X_{j+n-1}^2
 \ge k^{-6/5}e^{2n}
 \ge (nM_n)^{-6/5}e^{2n}.
\]
Since \(M_n\le2e^n\) for \(n\ge1\),
\[
 \log\|T^n x\|^2
 \ge \frac45n-\frac65\log n-\frac65\log2.
\]
Dividing by \(\sqrt{8n\log\log n}=o(n)\) proves the assertion.
\end{proof}

\section{Scope and verification}
The construction exploits infinite spatial support: at time \(n\), the norm
simultaneously samples exponentially many remote length-\(n\) blocks.  It
does not contradict Lemma~7.15 for any fixed basis vector.  The proof uses
only an exact norm identity and the first Borel--Cantelli lemma.  The
accompanying script checks the numerical inequalities used above for a
range of \(n\); those checks are transcription aids, not a substitute for
the analytic proof.

A bounded arXiv search through 21 August 2026 using the paper id, title,
authors, and the exact conjecture found no later paper recording this
counterexample.  This is not an exhaustive literature review.

\begin{thebibliography}{9}
\bibitem{CFZ} G.~Cheng, X.~Fang, and S.~Zhu,
\emph{Random weighted shifts}, arXiv:1811.05761.
\end{thebibliography}
\end{document}
